\chapter{Arithmetic FLE: Source Scope and Status Corrections}
\label{v4-arith:fle-audit}

The critical-level fundamental local equivalence (FLE) and its quantum
and small variants are theorems in the geometric Langlands programme
over complex curves. The arithmetic branch of this volume studies
\(\mathcal{V}^{\mathrm{prim}}\), the Hochschild trace class of the
identity on the global arithmetic Iwahori--Hecke category of
\(GL_{2}\), and its conjectural companions F1, F2, F3. This chapter
records the scope boundaries of the geometric FLE sources, explains
why F1/F2/F3 require a distinct mixed-characteristic framework, and
corrects the status of claims that previously conflated the two
settings. No \texttt{RealizedHere} decorator is installed here;
arithmetic-FLE realization remains outside the registry pending the
obligations recorded in Section~\ref{v4-arith:fle-audit:open-obligations}.

\section{Source Theorems}
\label{v4-arith:fle-audit:sources}

\begin{theorem}[geometric critical-level FLE]
\label{v4-arith:fle-audit:thm:critical-FLE-source}
\ClaimStatusProvedElsewhere. The 2024 geometric-Langlands proof
installs the critical-level fundamental local equivalence in the
geometric setting over a smooth projective curve in characteristic
\(0\):
\[
\mathrm{KL}(G)_{\mathrm{crit}}\;\simeq\;
\mathrm{IndCoh}^{*}\bigl(\mathrm{Op}^{\mathrm{mf}}_{\check G}\bigr),
\]
with compatibility with Kac--Moody localization. Its source is
Arinkin--Beraldo--Campbell--Chen--Faergeman--Gaitsgory--Lin--Raskin--
Rozenblyum, \emph{Proof of the geometric Langlands conjecture II:
Kac--Moody localization and the FLE}, arXiv:2405.03648.
\end{theorem}

\begin{theorem}[quantum and small FLE sources]
\label{v4-arith:fle-audit:thm:quantum-FLE-source}
\ClaimStatusProvedElsewhere. The quantum and small FLE layer consists
of local geometric and metaplectic/quantum statements over a complex
curve. Campbell--Dhillon--Raskin's \emph{Fundamental local
equivalences in quantum geometric Langlands} (arXiv:1907.03204) and
Gaitsgory--Lysenko's small-FLE work (arXiv:1903.02279) establish
equivalences whose source categories are Kazhdan--Lusztig or
metaplectic representations and whose target categories are quantum
groups or root-of-unity modules over a formal disk. The arithmetic
analogue of these equivalences is an open problem at each finite prime.
\end{theorem}

\begin{remark}[scope of the geometric FLE sources]
\label{v4-arith:fle-audit:rem:source-boundary}
The FLE sources operate within the geometric Langlands, Kac--Moody,
Whittaker, and root-of-unity quantum-group framework over complex
curves. A local categorical \(p\)-adic Langlands equivalence over
\(\mathbb{Q}_{p}\), a globalization over
\(\overline{\mathrm{Spec}\,\mathbb{Z}}\), and an Emerton--Gee full
Weil/\((\varphi,\Gamma)\) target comparison each require additional
technology: Fargues--Scholze's geometrization, prismatic or
pro-\'etale methods, and the Emerton--Gee eigenvariety construction.
\end{remark}

\section{Scope Distinction: Geometric FLE vs F1/F2/F3}
\label{v4-arith:fle-audit:non-identification}

\begin{proposition}[F1/F2/F3 as a distinct mixed-characteristic programme]
\label{v4-arith:fle-audit:prop:not-FLE}
\ClaimStatusNeedsVerification. The arithmetic branch's F1 (local
categorical Langlands comparison for \(GL_{2}/\mathbb{Q}_{p}\) against
the Emerton--Gee stack), F2 (mixed-characteristic \(E_{2}\) Hecke
enhancement and restricted tensor product), and F3 (Ben-Zvi--Nadler
Hochschild-trace comparison against the full Weil/\((\varphi,\Gamma)\)
target) constitute a mixed-characteristic local-global Langlands and
Hochschild-trace programme whose source and target categories are
disjoint from those of the geometric FLE.
\end{proposition}

\begin{proof}
The geometric FLE (Theorem~\ref{v4-arith:fle-audit:thm:critical-FLE-source})
pairs critical-level Kac--Moody or Whittaker/factorizable categories
as source with oper or quantum-group categories over a smooth
projective complex curve as target. F1 pairs the spherically-finite
automorphic category for \(GL_{2}/\mathbb{Q}\) as source with
ind-coherent sheaves on the Emerton--Gee stack as target. F2 requires
a mixed-characteristic \(E_{2}\) Hecke enhancement and restricted
tensor product. F3 requires a Ben-Zvi--Nadler Hochschild-trace
comparison against the full Weil/\((\varphi,\Gamma)\) arithmetic
target. The targets of F1, F2, F3 are Emerton--Gee and
Weil/\((\varphi,\Gamma)\) stacks; the inputs at finite primes are
adelic Hecke categories, not Beilinson--Drinfeld chiral algebras over
a complex curve.
\end{proof}

\begin{remark}[finite-prime categorical structure]
\label{v4-arith:fle-audit:rem:finite-prime-replacement}
At a finite prime \(p\), the relevant categorical structure is the
category \(\mathrm{Coh}(\mathrm{Bun}_{G})\) of coherent sheaves on
the moduli stack of \(G\)-bundles on the Fargues--Fontaine curve,
equipped with the Fargues--Scholze spectral action
(arXiv:2102.13459), together with Hecke correspondences, diamonds,
and \((\varphi,\Gamma)\)-modules.
The archimedean fibre carries classical Kac--Moody and chiral
intuition, providing the motivating analogy; the finite-prime
framework is governed by Fargues--Scholze geometrization and
Emerton--Gee eigenvarieties.
\end{remark}

\section{Scope Corrections Across Source Categories}
\label{v4-arith:fle-audit:ledger}

\begin{remark}[local/global type distinction]
\label{v4-arith:fle-audit:rem:local-global-type}
F1, F2, and F3 are local-global Langlands and Hochschild-trace
conjectures. Their source is the automorphic category for
\(GL_{2}/\mathbb{Q}\) or its local \(p\)-adic analogue; their target
is the Emerton--Gee stack or the full Weil/\((\varphi,\Gamma)\)
coherent sheaf theory. These source and target types differ from those
of the geometric FLE.
\end{remark}

\begin{remark}[finite-prime chiral structure]
\label{v4-arith:fle-audit:rem:chiral-at-finite-prime}
At a finite prime, the relevant structure is the
Beilinson--Drinfeld factorization algebra and the Hecke
correspondence on \(\mathrm{Bun}_{G}\), governed by Fargues--Scholze
geometrization. The archimedean fibre carries Kac--Moody and chiral
intuitions as a motivating analogy.
\end{remark}

\begin{remark}[quantum-group and root-of-unity source scope]
\label{v4-arith:fle-audit:rem:qgroup-scope}
The root-of-unity and small-FLE sources (Campbell--Dhillon--Raskin,
Gaitsgory--Lysenko) are geometric and quantum: they operate over a
complex curve with a quantum parameter \(q\) a root of unity or with
a metaplectic covering. The parameter \(p^{-s}\) appearing in local
Hecke eigenvalues at a prime \(p\) is a spectral parameter, distinct
from a root-of-unity quantum-group parameter.
\end{remark}

\begin{remark}[function-field transport and number-field scope]
\label{v4-arith:fle-audit:rem:function-field-transport}
Raskin's globalization operates over a function field. The transport
to a number field requires: (a) a local categorical comparison F1 at
each prime, and (b) a prismatic or BC-diamond-glue transition across
primes. Both remain open.
\end{remark}

\begin{remark}[mixed-characteristic substrate]
\label{v4-arith:fle-audit:rem:mixed-char-substrate}
The arithmetic compactification
\(\overline{\mathrm{Spec}\,\mathbb{Z}}\) carries an archimedean place
and mixed-characteristic fibres. A curve in characteristic \(0\)
carries neither. The arithmetic programme requires mixed-characteristic
methods (prismatic cohomology, perfectoid geometry) at each finite
prime and complex-analytic methods at the archimedean place.
\end{remark}

\begin{remark}[HZ-IV realization status for arithmetic FLE]
\label{v4-arith:fle-audit:rem:hz-iv-status}
Arithmetic FLE has no source-disjoint derivation/verification pair in
the realization registry. A \texttt{RealizedHere} decorator for
arithmetic FLE requires first closing the obligations in
Section~\ref{v4-arith:fle-audit:open-obligations}.
\end{remark}

\begin{remark}[AGKRRV 2024 source and target]
\label{v4-arith:fle-audit:rem:agkrrv-target}
Arinkin--Beraldo--Campbell--Chen--Faergeman--Gaitsgory--Lin--Raskin--%
Rozenblyum (arXiv:2405.03648) establishes the critical FLE between
\(\mathrm{KL}(G)_{\mathrm{crit}}\) and
\(\mathrm{IndCoh}^{*}(\mathrm{Op}^{\mathrm{mf}}_{\check G})\): source
is critical-level Kac--Moody, target is monodromic opers on a complex
curve. The Emerton--Gee stack is a mixed-characteristic eigenvariety;
its objects are \((\varphi,\Gamma)\)-modules, not monodromic opers.
\end{remark}

\begin{remark}[F3/BZN status against full Weil target]
\label{v4-arith:fle-audit:rem:f3-bzn}
F3 asks for a Ben-Zvi--Nadler Hochschild-trace identification between
the automorphic side and coherent sheaves on the Weil/\((\varphi,\Gamma)\)
stack. The inertia-only loop object
\(B\widehat{\mathbb{Z}}\) captures only the tame/procyclic part of
inertia; the full Weil group retains wild inertia and Frobenius data.
F3 with the full Weil target is a conjecture.
\end{remark}

\begin{remark}[spectral parameter vs quantum-group parameter]
\label{v4-arith:fle-audit:rem:q-vs-ps}
The Hecke eigenvalue at a prime \(p\) is recorded by the spectral
parameter \(p^{-s}\) or equivalently by the Satake image in
\(\mathrm{Rep}(\check G)\). The quantum-group deformation parameter
\(q\) appearing in the quantum FLE is a root of unity or a formal
parameter; identifying \(q = p^{-s}\) requires an explicit source-secure
bridge; this bridge is the sixth obligation in
Section~\ref{v4-arith:fle-audit:open-obligations}.
\end{remark}

\begin{remark}[Fargues--Scholze, Emerton--Gee, and Dotto--Emerton--Gee]
\label{v4-arith:fle-audit:rem:fs-eg-deg}
Fargues--Scholze (arXiv:2102.13459) provides the spectral action on
\(\mathrm{Coh}(\mathrm{Bun}_{G})\) via the Fargues--Fontaine curve.
Emerton--Gee (arXiv:1908.07185) constructs the stack of
\((\varphi,\Gamma)\)-modules. Dotto--Emerton--Gee establishes full
faithfulness of a functor to mod-\(p\) representations; an
equivalence of \(\infty\)-categories between the automorphic and the
Emerton--Gee coherent side is the content of F1.
\end{remark}

\begin{remark}[F1/F2/F3 claim-status correction]
\label{v4-arith:fle-audit:rem:f123-status}
F1 carries status \emph{Theorem-waiting}: the local categorical
comparison is supported by Fargues--Scholze (arXiv:2102.13459) and
Emerton--Gee (arXiv:1908.07185) and requires the adelic-descent
and BC-diamond-glue steps. F2 carries status
\emph{Theorem-waiting}: the mixed-characteristic \(E_{2}\) Hecke
enhancement and restricted tensor product are formally accessible but
unproved. F3 carries status \emph{Conjecture}: the full
Weil/\((\varphi,\Gamma)\) target comparison is beyond current
techniques. Chapters previously marking F1, F2, F3 with premature
theorem-closure have been downgraded accordingly.
\end{remark}

\begin{remark}[local Satake/Yangian normalization]
\label{v4-arith:fle-audit:rem:satake-yangian}
The local Satake isomorphism identifies the spherical Hecke algebra
with the representation ring of \(\check G\). An arithmetic Yangian
structure, if it exists, must be grounded in a source-secure
comparison between the Satake filtration and an explicit Yangian
coproduct; this identification remains open.
\end{remark}

\begin{remark}[BC-diamond-glue decomposition]
\label{v4-arith:fle-audit:rem:bc-diamond-glue}
The BC-diamond-glue problem decomposes into three parts: (a) the
local BZCHN/Emerton--Gee comparison at each prime \(p\), (b)
cross-prime prismatic transition maps, and (c) global
Bernstein-centre compatibility. Each part is open.
\end{remark}

\begin{remark}[wild-prime correction]
\label{v4-arith:fle-audit:rem:wild-prime}
At \(p = 2\) and \(p = 3\), the source for the local comparison
carries ramification and wild inertia data not present in the tame
case. A conjectural correction accounts for this; the precise form
depends on the source used for the local comparison.
\end{remark}

\section{Corrected Claim Status}
\label{v4-arith:fle-audit:status-matrix}

\begin{theorem}[geometric FLE: scope and status]
\label{v4-arith:fle-audit:thm:geometric-fle-status}
\ClaimStatusProvedElsewhere. The critical-level FLE
\(\mathrm{KL}(G)_{\mathrm{crit}} \simeq
\mathrm{IndCoh}^{*}(\mathrm{Op}^{\mathrm{mf}}_{\check G})\) and the
quantum/small FLE equivalences (Campbell--Dhillon--Raskin,
Gaitsgory--Lysenko) are proved in the geometric Langlands setting over
a complex curve. Their source is Campbell--Dhillon--Raskin
(arXiv:1907.03204), Gaitsgory--Lysenko (arXiv:1903.02279), and
Arinkin--Beraldo--Campbell--Chen--Faergeman--Gaitsgory--Lin--Raskin--%
Rozenblyum (arXiv:2405.03648). These results hold over smooth
projective curves in characteristic \(0\).
\end{theorem}

\begin{remark}[archimedean chiral structure and finite-prime scope]
\label{v4-arith:fle-audit:rem:chiral-scope}
The Kac--Moody and chiral algebra structures provide the motivating
analogy at the archimedean place. At finite primes, the geometric
structures are Beilinson--Drinfeld factorization algebras,
Hecke correspondences, and the Fargues--Scholze spectral action; the
finite-prime analogue of the critical-level FLE is the content of F1.
\end{remark}

\begin{remark}[corrected status of F1, F2, F3, and BC-diamond-glue]
\label{v4-arith:fle-audit:rem:corrected-status}
F1 (local categorical Langlands comparison for \(GL_{2}/\mathbb{Q}_{p}\)
via Emerton--Gee, with Fargues--Scholze spectral action as input)
carries status \emph{Theorem-waiting}: the comparison requires
adelic-descent and BC-diamond-glue steps.
F2 (mixed-characteristic \(E_{2}\) Hecke enhancement and restricted
tensor product for the adelic assembly) carries status
\emph{Theorem-waiting}: the \(\infty\)-categorical Tate restricted
tensor product is accessible in principle but unproved.
F3 (Ben-Zvi--Nadler Hochschild-trace identification against the full
Weil/\((\varphi,\Gamma)\) target, with the full Weil group including
wild inertia and Frobenius data) carries status \emph{Conjecture}.
BC-diamond-glue decomposes into local BZCHN/Emerton--Gee
comparison, cross-prime prismatic transition maps, and global
Bernstein-centre compatibility; each component is open.
The wild-prime correction at \(p = 2, 3\) is a source-sensitive
conjecture; the precise form depends on the local comparison chosen.
\end{remark}

\section{Open Obligations}
\label{v4-arith:fle-audit:open-obligations}

\begin{enumerate}
\item Formulate the mixed-characteristic local categorical
equivalence F1: construct a comparison between the spherically-finite
automorphic \(\infty\)-category for \(GL_{2}/\mathbb{Q}_{p}\) and
ind-coherent sheaves on the Emerton--Gee stack, with the
Fargues--Scholze spectral action (arXiv:2102.13459) as input.
\item Extend the loop-space comparison from the inertia-only object
\(B\widehat{\mathbb{Z}}\) to the full Weil group (with wild inertia
and Frobenius data) and prove the flat pullback statement in the full
Weil/\((\varphi,\Gamma)\) setting.
\item Prove profinite formal-loop continuity for
\(\mathrm{Map}(B\widehat{\mathbb{Z}}, -)\) in the formal Noetherian
pro-\'etale setting, supplying the foundational input for the
inertia-only portion of F3.
\item Prove the mixed-characteristic \(E_{2}\) Hecke enhancement and
the \(\infty\)-categorical Tate restricted tensor product required for
F2.
\item Prove BC-diamond-glue in its three components: (a) the local
BZCHN/Emerton--Gee comparison, (b) cross-prime prismatic transition
maps, and (c) global Bernstein-centre compatibility.
\item Construct a source-secure bridge between the Satake spectral
parameter \(p^{-s}\) and the quantum-group deformation parameter
\(q\) of the quantum FLE, grounded in an explicit comparison of the
Satake filtration and the Yangian coproduct.
\item Establish a source-disjoint derivation/verification pair for
arithmetic FLE so that a \texttt{RealizedHere} decorator can be
registered in the HZ-IV realization registry.
\end{enumerate}

\begin{remark}[source role of the geometric FLE results]
\label{v4-arith:fle-audit:rem:integration-rule}
The geometric FLE (Theorems~\ref{v4-arith:fle-audit:thm:critical-FLE-source}
and~\ref{v4-arith:fle-audit:thm:quantum-FLE-source}), quantum
geometric Langlands, and Raskin's globalization provide motivating
analogies and local-model inputs for F1, F2, F3. Each such use
requires an explicit source-secure bridge; the bridges are the
obligations listed above.
\end{remark}
